\chapter*{\RL{ملخص}}
\addcontentsline{toc}{chapter}{Arabic Abstract}

\begin{RLtext}

\noindent يعرض هذا التقرير الأعمال المنجزة خلال مشروع نهاية الدراسة بوحدة \LR{BCG X}، الذراع التقنية للتصميم والهندسة التابعة لشركة \LR{Boston Consulting Group (BCG)}. يتناول المشروع منتَج \LR{Energy Transition Optimizer}، وهو أداة داخلية للتقييم الأولي لجدوى مشاريع \LR{Power-to-X}، طُوِّرت داخليًا ثم كُيِّفت لعدة عملاء. تقوم الأداة على محرّك تحسين حتمي: يُوصَف كل مشروع بطوبولوجيا (رسم بياني مُصنَّف من الأصول وعُقَد الحفظ والأسواق) وبمئات المعطيات التقنية والمالية، ينطلق منها المحرّك لمحاكاة التشغيل ساعةً بساعة وحساب مؤشرات كـ \LR{NPV} و\LR{IRR} والكلفة المُسوّاة للمنتَج \LR{(LCOX)}.

\noindent رغم قوته، ظل هذا المنتج حكرًا على الخبراء وقائمًا على منصة هشّة، فوجب تعزيزه وجعله في متناول غير المختصين. أُنجز العمل ضمن فريق، فحدّث المنصة وأضاف طبقة محادثة متعددة الوكلاء محكومة: يوجِّه وكيلٌ موجِّه كل طلب إلى وكلاء متخصصين (استيعاب البيانات، وتصميم الطوبولوجيا، واستكشاف السيناريوهات، وتفسير النتائج)، بينما يظل أي إجراء يُعدِّل حالة المشروع اقتراحًا مُنمَّطًا يخضع للتحقق ثم للمصادقة البشرية، ليبقى المحرّك الحتمي المرجع الوحيد للقيم العددية.

\noindent وتكتمل هذه الطبقة بثلاثة محاور: وكيل تفسير النتائج، وهو وكيل فرعي للقراءة فقط يشرح بلغة طبيعية اقتصاديات عمليات المحاكاة مع إبقاء كل رقم قابلًا للتتبّع حتى نتيجة حتمية دون أي قيمة مُختلَقة؛ وطبقة التمثيل البصري المرافِقة (مخططات تبعثر قابلة للتفسير عند الطلب وجدول تدفقات نقدية مخصومة قابل للتصفّح)؛ ومنظومة تقييم بنموذج لغوي حَكَم \LR{(LLM-as-a-judge)} تقيس الارتكاز الوقائعي والسلامة المالية للتفسيرات. وبذلك يجمع النظام بين قابلية إعادة الإنتاج والتدقيق وسهولة الوصول عبر واجهة مدعومة بالذكاء الاصطناعي.

\end{RLtext}

\noindent\rule[1pt]{\textwidth}{0.5pt}

\begin{RLtext}

\noindent{\textbf{الكلمات المفتاحية :}}
\LR{Power-to-X}، التقييم الأولي للجدوى، التحسين الحتمي، الأنظمة متعددة الوكلاء، النماذج اللغوية الكبيرة، تفسير النتائج، التقييم بنموذج لغوي حَكَم، المصادقة البشرية.
\end{RLtext}

\vspace{0.3cm}
\noindent\rule[1pt]{\textwidth}{0.5pt}
